\section{Yann LeCun 与 JEPA 框架}

谈到 Yann LeCun，谢赛宁的描述非常生动：网上的 LeCun 是“斗士”，现实中的 LeCun 则是“很好很好的人”。这种反差本身并不罕见，但在 LeCun 身上尤其有趣，因为他始终同时扮演科学共同体中的争论者与组织中的稳定者。一方面，他在公开场合坚定捍卫自己对 LLM 局限性的判断，不会为了公司叙事而轻易转口；另一方面，他又有极强的个人兴趣和生活感，造模型飞机、做天文摄影、听电子乐与 Jazz、开帆船，这些爱好让他看起来更像一个始终处于青春期延长期的人。何凯明那句 \textit{``LeCun 是一个 16 岁青春期一直延续到 65 岁的人''}，并不是调侃，而是对一种旺盛好奇心和玩心的精确概括。

\begin{knowledgebox}{LeCun 的人格结构}
谢赛宁看到的 LeCun，并不是单一的“图灵奖得主”形象，而是一个持续动手、持续好奇、持续公开争辩的人。模型飞机、天文、Jazz、帆船这些兴趣的共同点，在于都需要长期投入、技术感与审美感并存，也都要求人既尊重现实约束，又保留探索未知的兴奋。
\end{knowledgebox}

这种人格结构和他的管理方式是一致的。谢赛宁说 LeCun 管团队像开帆船：绝大多数时间，你需要充分信任每个人在自己的位置上发挥；但一旦看到风向、海流或船体姿态有偏差，就要尽早校正，而不是等问题扩大后再强硬拉回。对研究组织来说，这是一种很成熟的管理哲学。因为真正有创造力的人不能被过度 micromanage，可若完全放任，又容易让团队在局部兴趣中分散。帆船管理法在这里的核心，是让组织保持方向感，同时不过早扼杀每个人独立发现问题的能力。

JEPA 正是在这样的背景下被理解的。谢赛宁特别强调，JEPA 不是一个模型，而是一个非常广阔的海洋。它不是“某篇论文里的一个模块组合”，而是一套关于表征、预测与规划的总体框架。也因此，他把自己理解 JEPA 的过程总结成三个阶段：质疑 JEPA，理解 JEPA，成为 JEPA。这种说法听起来像玩笑，却准确抓住了很多前沿框架的认知路径。最初你会觉得它太抽象、太宽泛、像一个几乎什么都能往里装的概念；但当你反复看 LeCun 的 talk、把它放回世界模型语境中思考，就会逐渐意识到，JEPA 的广阔不是含糊，而是因为它确实试图覆盖从 world understanding 到 prediction 再到 planning 的整条链路。

\begin{importantbox}{JEPA 的三阶段认知}
\begin{tabularx}{\textwidth}{>{\raggedright\arraybackslash}p{0.2\textwidth} Y}
\toprule
阶段 & 典型状态 \\
\midrule
质疑 JEPA & 觉得它过于抽象，不像一个“立刻可跑”的 recipe，难以和现成 benchmark 对齐。 \\
理解 JEPA & 开始意识到它在描述一种更高层的学习目标：不是拟合像素或 token，而是学习对世界状态有预测价值的抽象表示。 \\
成为 JEPA & 不再把它当成某个单独方法，而是把表征、预测、规划与行动统一放入同一框架下思考。 \\
\bottomrule
\end{tabularx}
\end{importantbox}

这种框架思维也支撑了 LeCun 对 LLM 的长期批评。谢赛宁引用那句极具戏剧感的话：\textit{``LLM 终将凋零...老兵不死，终将凋零''}。这并不是说语言模型会突然失效，而是说它们不会成为构建通用智能系统的最终基石。原因很清楚：如果一个系统主要通过文本序列学习，它就在很大程度上被限制在“他人已经压缩好的世界描述”里，而没有充分建立起面对现实的自有世界模型。JEPA 的雄心，正是要绕开这种局限，去学习更接近世界结构本身的表示。

\begin{warningbox}{把 LLM 当作终局，是一种阶段性幻觉}
LeCun 和谢赛宁并不否认 LLM 的巨大价值，但他们都反对把今天的语言模型能力直接外推成终局。语言是极强的抽象工具，却不是全部现实。若一个系统主要依赖语言对齐世界，它就仍然像拄着拐杖前行。拐杖可以帮你站起来，却未必足以让你跑进真实世界。
\end{warningbox}

LeCun 在谢赛宁这里之所以有说服力，还因为他始终保留科学家式的 integrity。访谈中一个很重要的观察是：即使所在公司、所处风向、资本市场的偏好都在推动“向 LLM 靠拢”的叙事，LeCun 也不会因为组织需要而公开修改自己对技术路线的根本判断。这种坚持并不总让人舒服，却对一个研究共同体极其重要。因为如果最核心的问题判断都可以随着风口随意改写，那么所谓研究方向就会退化成营销话术的延长线。

谢赛宁一再观看 LeCun talk 的经历，也说明 JEPA 并不是那种“听完就懂”的框架。它要求研究者逐渐建立更高层的抽象习惯：不要急着问它对应哪一个榜单、哪一个 loss、哪一个模块，而要先问它到底在刻画什么学习目标。正因为这种理解过程很慢，JEPA 才容易在早期被误判成空泛；也正因为它足够宽，才可能在未来容纳更完整的世界理解与规划体系。

谢赛宁还提到，自己把 LeCun 的 talk 看了十到二十遍，每次都有新收获。这句话很能说明 JEPA 的理解门槛。真正宽的框架，不会在第一次听时就被完全吸收；它需要你带着不同阶段的经验反复返回。某种意义上，这与他对研究的理解高度一致：好的理论不会只给你一个结论，而会在你能力变化后继续提供新的可见度。

最后还有一个很温暖的比喻。谢赛宁说，\textit{``Yann 是一个巨大的电池，他赶着我，我希望把电力输送下去''}。这句话说明 LeCun 在他心中并不仅是思想导师，也是能量源。一个真正强的学术领袖，不只生产观点，还能把身边人的势能整体抬高，让他们愿意往更难的问题上持续投注。

\subsection{本章小结}

在谢赛宁眼里，LeCun 之所以重要，不只是因为其历史地位，而是因为他把人格、管理和方法框架统一起来。JEPA 的广阔、对 LLM 终局论的警惕、帆船式管理与“巨大电池”式的带动能力，共同解释了为什么谢赛宁愿意把下一阶段押在与他同行之上。

